\documentclass[../../../include/open-logic-section]{subfiles}

\begin{document}

\olfileid[pt]{sth}{choice}{woproblem}
\olsection{O Problema da Boa Ordenação}

Evidentemente que muito depende de aceitarmos ou não a Boa Ordenação. Mas
a discussão deste princípio tendeu a concentrar-se num equivalente,
o Axioma da Escolha. Voltemos então a atenção para esse axioma
(e provemos a equivalência). 

Em \citeyear{Cantor1883}, Cantor exprimiu o seu apoio ao Axioma
da Boa Ordenação, chamando-lhe ``uma lei do pensamento que me parece
fundamental, rica nas suas consequências e particularmente notável
pela sua generalidade'' (citado em \citeauthor{Potter2004}
\citeyear[p.~243]{Potter2004}). Mas Cantor acabou por convencer-se
de que o ``Axioma'' precisava de prova. A comunidade matemática
também. 

O problema foi ``resolvido'' por Zermelo em \citeyear{Zermelo1904}. Para
explicar a solução, precisamos de algumas definições. 

\begin{defn}
Uma função $f$ é uma \emph{função de escolha} sse $f(x) \in x$ para todo o $x \in \dom{f}$. Dizemos que $f$ é uma \emph{função de escolha para $A$} sse $f$ é uma função de escolha com $\dom{f} = A \setminus \{\emptyset\}$.
\end{defn}

Intuitivamente, para cada conjunto (não vazio) $x \in A$, uma função de escolha
para $A$ \emph{escolhe} um elemento particular, $f(x)$, de $x$. O
Axioma da Escolha é então:

\begin{axiom}[Escolha]
	Todo o conjunto tem uma função de escolha.
\end{axiom}

Zermelo mostrou que a Escolha implica a boa ordenação, e vice-versa:

\begin{thm}[em $\ZF$]\ollabel{thmwochoice}
A Boa Ordenação e a Escolha são equivalentes.
\end{thm}

\begin{proof}
\emph{Da esquerda para a direita.} Seja $A$ um conjunto de conjuntos. Então $\bigcup A$
existe pelo Axioma da União, e logo, pela Boa Ordenação, existe alguma relação
$<$ que bem ordena $\bigcup A$. Seja agora $f(x) = \text{o membro $<$-mínimo de }x$. Isto é uma função de escolha para $A$.

\emph{Da direita para a esquerda.} Fixe-se $A$. Pela Escolha, existe uma função de escolha,
$f$, para $\Pow{A} \setminus \{\emptyset\}$. Usando Recursão Transfinita, defina-se uma função:
\begin{align*}
	g(0) &= f(A)\\
	g(\alpha) &= 
		\begin{cases}
			\text{parar!{}} &\text{se }A = \funimage{g}{\alpha}\\
			f(A \setminus \funimage{g}{\alpha}) & \text{caso contrário}\\	
		\end{cases}
\end{align*}
A indicação ``parar!'' é apenas uma abreviatura para o que seria
de outro modo uma definição mais longa. Ou seja, quando $A =
\funimage{g}{\alpha}$ pela primeira vez, seja $g(\delta) = A$ para todo o
$\delta \leq \alpha$. Agora, em primeiro lugar, só podemos ter a certeza de que isto define um \emph{termo} (ver as observações após \olref[sth][spine][recursion]{transrecursionschema}); mas mostraremos que de fato temos uma função.

Como $f$ é função de escolha, para cada $\alpha$ (quando definido) temos $g(\alpha) =
f(A \setminus \funimage{g}{\alpha}) \in A \setminus
\funimage{g}{\alpha}$; isto é, $g(\alpha) \notin \funimage{g}{\alpha}$.
Logo, se $g(\alpha) = g(\beta)$ então $g(\beta) \notin
\funimage{g}{\alpha}$, isto é, $\beta \notin \alpha$, e analogamente
$\alpha \notin \beta$. Logo $\alpha = \beta$, por Tricotomia. Assim, $g$ é
!!{injective}.

Observe-se a seguir que de fato paramos!{}, isto é, que existe algum ordinal mínimo $\alpha$ tal que $A = g[\alpha]$. De fato, suponhamos o contrário; então, como $g$ é !!{injective}, teríamos $\cardless{\alpha}{\Pow{A} \setminus \{\emptyset\}}$ para
todo o ordinal $\alpha$, contradizendo \olref[hartogs]{HartogsLemma}. Daí também $\ran{g} = A$.

Reunindo estes fatos, $g$ é !!a{bijection} de algum ordinal para $A$. Agora, $g$ pode ser usada para bem ordenar $A$.
\end{proof}

Assim, a Boa Ordenação e a Escolha estão de pé ou caem juntas. Mas a questão
permanece: estão de pé ou caem?

\end{document}
